Thus, unless P${}={}$NP, neither parameter can be dropped and the dependence
on $\kappa$ cannot be polylogarithmic; under ETH the dependence on $p$ is
exponential even for $\kappa\le2$; and under rETH the exponent of $\kappa$
cannot be $o(p)$, already for integer box quadratics, whereas
Theorem~\ref{thm:intro-main}(b) has exponent $p/2+O(1)$. Because
$\bar\kappa\le\kappa$, the same statements hold for $\bar\kappa$. In the
value-oracle model of part~(d), CT attains the exponent $p/2$ up to a factor
$(C\sqrt p\log(p+2))^p$ and the number of stages
(Remark~\ref{lim:rem:oracle}). Section~\ref{sec:limits} also shows that
growth towards a continuum of minimizers bounds neither the grids produced by
filtering nor, independently of the accuracy, the number of value and
derivative queries that a correct deterministic algorithm needs; that a path of local affine equalities makes the problem
NP-hard at bag size three and $\kappa=1$; and that matching separator
moments of any finite order gives no local error bound in terms of negative
curvature.
As far as we know, the expanding-box chain, the unique-minimizer hardness
with explicit growth (a refinement of \cite[Remark~2]{DelPiaKhajavirad2026}),
the rETH bound and the moment obstruction are new. The other statements
adapt standard constructions, among them the running-sum encoding of Subset
Sum of Bienstock and Mu\~noz \cite{BienstockMunoz2018} (see also
\cite[Example~1.1]{CifuentesParrilo2016}).
%%%
\paragraph{Conditional recourse.}
The condition number is built from diagonal curvature, so a stiff convex
term inflates it even when the subproblem it defines is easy. We therefore
also allow \emph{recourse}: exact minimization over some variables $y$ for
fixed values $v$ of the others. Partial minimization over a box preserves
upper coordinate curvatures and growth, so the certificate and the growth
analysis apply to the value function $V(v)=\min_yF(v,y)$ whenever $V$ can be
evaluated exactly (Theorem~\ref{thm:valuefn}). Two classes have exact recourse with
certificates. For private convex quadratic blocks, the optimal responses are
piecewise affine and certify a reduction of the coordinate curvatures of $V$
that is the largest possible for the block (Theorem~\ref{thm:cv} and
Proposition~\ref{prop:vf-curv}). If a quadratic $F$ is concave in each
recourse variable and submodular in them after sign changes, each recourse
problem is one minimum cut with a flow certificate, so only the remaining
variables are gridded (Theorems~\ref{thm:cr-oracle} and~\ref{thm:cr-search}).
For balanced box quadratics the grid problem itself is a minimum cut, so no
tree decomposition is needed: under point growth, certified approximation
takes $(n\kappa(I+q+1))^{O(1)}$ bit operations (Theorem~\ref{thm:balanced}).
Recourse does not reduce the parameter to the ratio of negative curvature to
growth in general (Proposition~\ref{prop:cv-limit}). Parametric quadratic
programming, cut representations of submodular quadratic functions and
threshold encodings of ordered labels are classical; what is new is the
curvature and growth accounting that keeps them inside the certificate.
%%%
\paragraph{Coupling constraints.}
The certificate rests on rounding each coordinate independently, which
breaks under coupling constraints. Totally unimodular (TU) coupling
constraints with mesh-aligned data admit exactly feasible correlated
rounding, which restores the certificate and, for quadratics, exact output
(Section~\ref{sec:constraints}). Under set growth with finitely many optimal
values per coordinate, the algorithm is polynomial for fixed width when the
condition number $\kappa_c$ (the ratio of a curvature bound for the
continuous variables to the set-growth constant), the box widths measured in
mesh units ($s/\eta$) and the number $r$ of optimal values per coordinate are
polynomially bounded (Theorem~\ref{thm:tu-approx}). It is not
fixed-parameter tractable in $(p,\kappa_c)$, because its meshes are uniform
(Section~\ref{sec:tu-limits}), and its pseudopolynomial first grid cannot be
removed unless P${}={}$NP (Proposition~\ref{lim:prop:constraints}). Bienstock
and Mu\~noz \cite{BienstockMunoz2018} allow arbitrary polynomial constraints
but satisfy them only within a tolerance; our class is narrower but exactly
feasible (Remark~\ref{rem:tu-bm}).
%%%
\paragraph{Several minimizers.}
One graded grid can need exponential time already with two isolated
minimizers (Proposition~\ref{prop:twocenters}). If the optimal set is finite,
unions of uniform cells keep $O(r\sqrt{n\kappa_S})$ nodes per coordinate,
where $r$ bounds the number of optimal values of a coordinate and
$\kappa_S=\max\{1,L/g_S\}$ uses the growth constant $g_S$ towards the optimal
set (Theorem~\ref{thm:cells}). For quadratics with nonpositive Hessian diagonal
on a mixed box, a dynamic program over the box endpoints describes the whole
optimal set as a constraint problem on the same tree decomposition. From it,
uniqueness, the dimension of the optimal set and, if it is finite, the number
of minimizers are computed in $2^{O(p)}\poly(I)$ bit operations
(Theorem~\ref{thm:endpointset} and Corollary~\ref{cor:facecsp}). This
combines Rosenberg's description of the optimal set of a multilinear function
as a union of faces \cite{Rosenberg1972} with zero-residual descriptions of
all optimal labelings of discrete problems
\cite{WainwrightJaakkolaWillsky2005,Werner2007}; the treatment of strictly
concave coordinates and mixed boxes and the factored form are new. For continuous box quadratics with a diagonal
Lagrangian certificate of optimality, which is classical
\cite{JeyakumarRubinovWu2006,LiWuQuan2015}, an exact minimizer and an exact
description of the optimal set are found in $f'(p,\kappa_S)\poly(I)$ bit
operations, for a computable function $f'$, without knowledge of $g_S$
(Theorem~\ref{thm:diagdiscovery}); we use that the same certificate holds at
every minimizer.
